\documentclass[10pt]{article}

\usepackage[margin=0.8in]{geometry}
\usepackage{booktabs}
\usepackage[hidelinks]{hyperref}

\title{SafeRL-Drive Phase-1 Surrogate Notes}
\author{Dheeraj Dhillon}
\date{\today}

\begin{document}
\maketitle

\section{Purpose}
This document is the detailed internal companion to \texttt{main.tex}. It is intended for
experiment bookkeeping, future Codex context, failed attempts, and decisions that would
be distracting in the portfolio report. Update it after each substantial Colab session.

\section{2026-07-21 diagnosis and Priority-0 correction}
The latest imported artifacts were:
\begin{itemize}
  \item IDM: \texttt{20260721\_023109\_idm\_baseline\_seed0};
  \item PPO: \texttt{20260721\_023556\_ppo\_mvp\_ppo\_seed0};
  \item SAC: \texttt{20260721\_032424\_sac\_mvp\_sac\_seed0}.
\end{itemize}

The best PPO checkpoint achieved 16\% success, 12\% collision, 54\% off-road,
18\% timeout, and 59.8\% mean route completion on the old 50-episode evaluation.
This is a meaningful improvement over the earlier 2\% PPO result, but road departure
remains the dominant failure. PPO training diagnostics were numerically healthy; its
validation reward peaked near 300,000 steps and later updates did not reliably improve
task success. The previous Stable-Baselines3 callback selected checkpoints by mean reward,
which was not aligned with the headline success metric.

SAC did produce complete artifacts, but the policy collapsed. Its best old evaluation had
0\% success, 14\% off-road, 86\% timeout, and 17.5\% route completion. The final model had
0\% success, 98\% timeout, and only 1.1\% route completion. The learned entropy coefficient
was approximately 0.014 at 150,000 steps, 0.211 at 300,000, 3.14 at 400,000, and 13.24 at
500,000. Critic loss rose from less than 1 around 150,000 steps to hundreds at 300,000 and
millions near 400,000. Actor-loss magnitude grew past 11,000. The synchronized optimizer,
entropy, and evaluation collapse indicates unstable value learning rather than a missing
output or a GPU-recording failure. The final policy mostly braked or used saturated actions,
making late timeout a low-immediate-cost failure mode.

Three measurement defects were found. First, MetaDrive's \texttt{info["velocity"]} is
already in km/h, but the evaluator multiplied it by 3.6; every legacy mean-speed result is
therefore inflated and must not be quoted. Second, repeated IDM evaluations over nominally
identical seeds produced 42\%, 30\%, and 50\% success because evaluation inherited random
traffic and did not seed every relevant random generator consistently. Third, seeds
1000--1049 were used for both checkpoint selection and final reporting, making them a
validation set rather than an untouched test set.

Priority 0 corrects these issues as follows:
\begin{itemize}
  \item training uses seeds beginning at 0 and keeps randomized traffic;
  \item deterministic validation uses seeds 1000--1049;
  \item deterministic held-out testing uses seeds 2000--2099 and runs only after the
        validation winner is frozen;
  \item the best checkpoint is ranked by success, then route completion, lower off-road,
        collision and timeout rates, and finally mean return;
  \item checkpoint normalization statistics are saved every 50,000 steps and beside the
        best model;
  \item training outcomes and PPO/SAC optimizer diagnostics are saved to
        \texttt{training\_diagnostics.json};
  \item success summaries include a 95\% Wilson interval;
  \item video recording explicitly sets the configured scenario seed immediately before
        reset and writes the seed and final outcome beside the MP4.
\end{itemize}

The old results remain useful diagnostic evidence, but are not final Phase-1 benchmark
numbers because their traffic realizations were not paired, their test seeds were reused
for model selection, and their speed field had the unit error.

The corrected local Priority-0 validation backfill used fixed traffic on seeds
1000--1019. PPO achieved 5\% success (95\% Wilson interval 0.9--23.6\%), 15\%
collision, 65\% off-road, 15\% timeout, and 47.1\% mean route completion. SAC
achieved 0\% success (95\% Wilson interval 0--16.1\%), 0\% collision, 20\%
off-road, 80\% timeout, and 14.2\% mean route completion. Corrected mean speeds
were 6.36 km/h for PPO and 1.25 km/h for SAC. These 20-episode validation results
confirm the qualitative diagnosis: PPO makes useful progress but usually leaves the
road, whereas SAC mostly stalls. They are diagnostic results, not held-out test claims.

The reproducibility gate then evaluated IDM twice in fresh worker environments on
seeds 1000--1009. All per-seed categorical outcomes matched exactly. Mean return
differed by 0.584 against an allowed 7.816, mean length by 0.5 against 11.044,
route completion by 0.00055 against 0.01, and mean speed by 0.00053 km/h against
0.254. The gate therefore passed. The first 10-episode pass had 40\% success and
is retained only as gate evidence, not as the final IDM benchmark.

\section{2026-07-21 newest-run audit and ground-up correction}
The later Google Drive import supersedes the artifacts listed above for diagnosis. The
newest completed learned runs are
\texttt{20260721\_090542\_ppo\_mvp\_ppo\_seed0} and
\texttt{20260721\_105717\_sac\_mvp\_sac\_seed0}; the newest IDM test is
\texttt{20260721\_181208\_idm\_baseline\_seed0}. Only these timestamps were used for the
second audit.

The frozen PPO winner reached 4\% success, 10\% collision, 75\% off-road, 11\%
timeout, and 32.5\% route completion on 100 test episodes. Its validation series peaked
at 10\% success at 350,000 steps and fell back to 0\% at 500,000. PPO's explained
variance and other optimizer diagnostics were mostly healthy. The agent was learning the
configured return, but that return was not a good enough proxy for safe completion. In
particular, successful episodes returned about 480 on average, while out-of-road, crash,
and timeout episodes still returned about 88, 103, and 240. Failure remained profitable.

The frozen SAC winner reached 0\% success, 7\% off-road, 93\% timeout, 15.5\% route
completion, and only 0.99 km/h mean speed. The failure is not missing evaluation output.
The critic loss grew from 0.17 at 210,000 steps to 61.9 at 310,000, 1,858 at 410,000,
and 5,817 at 500,000, with a peak above 12,000. At the same time the learned entropy
coefficient grew from roughly 0.01 to 1.36 and the actor-loss magnitude exceeded 1,000.
The best validation point was at 200,000 steps; later learning collapsed toward braking
and stalling. Timeout episodes still returned about 72 on average, so the replay learner
had both an attractive behavioral loophole and unstable value targets.

The ground-up code audit found the following causal defects:
\begin{enumerate}
  \item MetaDrive replaces only the terminal step reward on success or failure; it does
        not erase the dense progress accumulated earlier. Terminal penalties of 20 could
        not dominate hundreds of units of dense return.
  \item \texttt{use\_lateral\_reward} only scales positive progress near a lane edge; it
        is not a negative lane-centering cost. The prior documentation overstated it.
  \item No timeout penalty was configured, and \texttt{truncate\_as\_terminate} was
        false. Stable-Baselines3 therefore bootstrapped through horizon truncations. That
        is inappropriate here because this project explicitly counts timeout as failure,
        and it was especially damaging to SAC's replay targets.
  \item MetaDrive's state/lidar observation is already bounded in [0,1]. A second adaptive
        \texttt{VecNormalize} transform added nonstationarity without evidence of benefit.
  \item The old task exposed only 50 training scenarios on a five-block map with density
        0.1. It was simultaneously narrow for generalization and difficult enough that
        even native reference policies often failed.
  \item Earlier telemetry retained terminal actions but not full-rollout action behavior or
        the split between base reward and shaping penalties, making collapse harder to see.
\end{enumerate}

Native-policy sanity checks on validation seeds 1000--1019 separated simulator difficulty
from optimizer behavior:
\begin{center}
\begin{tabular}{lrr}
  \toprule
  Task & Expert success & IDM success \\
  \midrule
  Map 5, density 0.10 & 55\% & 35\% \\
  Map 3, density 0.00 & 100\% & 100\% \\
  Map 3, density 0.05 & 100\% & 100\% \\
  Map 4, density 0.05 & 95\% & 95\% \\
  Map 4, density 0.10 & 65\% & 50\% \\
  \bottomrule
\end{tabular}
\end{center}
This evidence justifies using map 3 at density 0.05 as the Phase-1 learning benchmark.
This is an intentional simplification, not an algorithmic breakthrough. Higher map and
traffic difficulty belongs in a later robustness phase after basic learning is established.

The new implementation uses 500 training scenarios, fixed spawn lanes, direct bounded
observations, terminal treatment of timeouts, success reward 100, failure penalties
50, and forward-progress weight 0.2. A wrapper penalizes speed below 5 km/h, lateral lane
deviation, horizon timeout, steering magnitude, and abrupt steering changes while logging the base reward and
each penalty. PPO now uses learning rate $10^{-4}$, 4,096-sample rollouts, batch size 256,
and CPU. SAC uses learning rate $10^{-4}$, a 500,000-step replay buffer, 20,000 learning
starts, batch size 512, GPU when available, and fixed entropy coefficient 0.05. Full action
distribution metrics are saved during training and evaluation.

Because seeds 2000--2099 informed this diagnosis, the untouched final test range moves to
3000--3099. Before another 500,000-step run, each algorithm must pass a 100,000-step
no-traffic control pilot on fixed validation seeds with at least 10\% success and 50\%
route completion. The SAC pilot also fails on a non-finite or final critic loss above
1,000. Pilot pointers are named \texttt{latest\_ppo\_pilot.txt} and
\texttt{latest\_sac\_pilot.txt}, so they cannot replace full-run pointers. If a gate
fails, stop; more timesteps are not the next experiment.

\section{Canonical experiment order}
\begin{enumerate}
  \item Run the smoke test; do not use it as a report result.
  \item Verify reproducible evaluation by running IDM twice on the same 10 validation seeds,
        requiring identical categorical outcomes and near-identical aggregate metrics.
  \item Require the native expert to achieve at least 80\% success on 10 validation seeds.
  \item Run the 100,000-step no-traffic PPO and SAC controls; require both learning gates.
  \item Train PPO and save its complete run directory only after its gate passes.
  \item Train SAC and save its complete run directory only after its gate passes.
  \item After both configurations are frozen, evaluate IDM and each validation winner
        once on the same 100 held-out test scenarios.
  \item Record one top-down rollout per learned agent.
  \item Generate the IDM--PPO--SAC comparison.
  \item Fill the report placeholders and compile \texttt{main.tex}.
\end{enumerate}

\section{Exact commands}
Run from the repository root:
\begin{verbatim}
python -m scripts.train --config configs/smoke_test.yaml
python -m scripts.evaluate_baseline --config configs/ppo_mvp.yaml \
  --policy expert --split validation --episodes 10 --prefix expert_sanity
python -m scripts.evaluate_baseline --config configs/ppo_mvp.yaml \
  --split validation --episodes 10 --prefix idm_repro --verify-repeat
python -m scripts.train --config configs/ppo_mvp.yaml --run-name ppo_control_pilot \
  experiment.latest_name=ppo_pilot train.total_timesteps=100000 \
  train.checkpoint_freq=25000 train.eval_freq=25000 validation.episodes=20 \
  metadrive.traffic_density=0.0 metadrive.random_traffic=false \
  validation.traffic_density=0.0
python -m scripts.train --config configs/sac_mvp.yaml --run-name sac_control_pilot \
  experiment.latest_name=sac_pilot train.total_timesteps=100000 \
  train.checkpoint_freq=25000 train.eval_freq=25000 validation.episodes=20 \
  metadrive.traffic_density=0.0 metadrive.random_traffic=false \
  validation.traffic_density=0.0
python -m scripts.train --config configs/ppo_mvp.yaml
python -m scripts.train --config configs/sac_mvp.yaml
python -m scripts.evaluate --run-dir <PPO_RUN> --model best \
  --split train --episodes 50 --prefix best_train_diagnostic
python -m scripts.evaluate --run-dir <SAC_RUN> --model best \
  --split train --episodes 50 --prefix best_train_diagnostic
python -m scripts.evaluate_baseline --config configs/ppo_mvp.yaml \
  --split test --episodes 100 --prefix idm_test
python -m scripts.evaluate --run-dir <PPO_RUN> --model best \
  --split test --episodes 100 --prefix best_test
python -m scripts.evaluate --run-dir <SAC_RUN> --model best \
  --split test --episodes 100 --prefix best_test
python -m scripts.record_video --run-dir <PPO_RUN> --model best \
  --seed <DISCLOSED_SUCCESSFUL_TEST_SEED> --steps 1000
python -m scripts.record_video --run-dir <SAC_RUN> --model best \
  --seed <DISCLOSED_SUCCESSFUL_TEST_SEED> --steps 1000
python -m scripts.compare_runs --phase1
latexmk -cd -pdf -interaction=nonstopmode -halt-on-error reports/main.tex
\end{verbatim}

\section{Phase-1 configuration record}
\begin{tabular}{lll}
  \toprule
  Setting & PPO & SAC \\
  \midrule
  Timesteps & 500,000 & 500,000 \\
  Policy & MLP & MLP \\
  Action space & Continuous & Continuous \\
  Vector environments & 4 subprocess & 1 dummy \\
  Observation normalization & No & No \\
  Map blocks / traffic density & 3 / 0.05 & 3 / 0.05 \\
  Training scenarios & 500 & 500 \\
  Training start seed & 0 & 0 \\
  Validation start seed & 1000 & 1000 \\
  Held-out test start seed & 3000 & 3000 \\
  Held-out test episodes & 100 & 100 \\
  \bottomrule
\end{tabular}

Exact algorithm hyperparameters live in \texttt{configs/ppo\_mvp.yaml} and
\texttt{configs/sac\_mvp.yaml}. Each run also contains the immutable
\texttt{resolved\_config.yaml} used for that run.

\section{Colab runtime metadata}
Record the following from each long run's \texttt{logs/train.log} and
\texttt{run\_metadata.json}:
\begin{itemize}
  \item Date and Colab runtime type: TODO
  \item Git commit: TODO
  \item Python, MetaDrive, Stable-Baselines3, and Gymnasium versions: TODO
  \item CUDA availability and GPU name: TODO
  \item Wall-clock start/end and interruptions: TODO
  \item Drive backup location: \texttt{/content/drive/MyDrive/SafeDrive}
\end{itemize}

\section{Observed results}
\begin{itemize}
  \item Initial IDM run (50 old evaluation episodes, original reward config): 42\% success,
        40\% collision, 16\% off-road, and 68.2\% mean route completion. Its legacy speed
        value was invalid because of the 3.6 multiplier bug.
  \item Initial PPO final model (500,000 steps, original reward config): 2\% success,
        40\% collision, 60\% off-road, and 27.1\% mean route completion. Its legacy speed
        was invalid. Its mean episode length was 195 steps versus IDM's 526, consistent
        with a faster policy terminating early through collisions or road departure.
  \item Latest PPO best-model diagnostic: 16\% success, 12\% collision, 54\% off-road,
        18\% timeout, and 59.8\% route completion.
  \item Latest SAC best-model diagnostic: 0\% success, 0\% collision, 14\% off-road,
        86\% timeout, and 17.5\% route completion.
  \item Priority-0 deterministic PPO validation backfill (20 episodes): 5\% success,
        15\% collision, 65\% off-road, 15\% timeout, and 47.1\% route completion.
  \item Priority-0 deterministic SAC validation backfill (20 episodes): 0\% success,
        0\% collision, 20\% off-road, 80\% timeout, and 14.2\% route completion.
  \item Initial interpretation: the original reward used weak terminal penalties and no
        lateral lane-centering factor. This could reward rapid forward progress before a
        failure. The revised config aligns terminal and lane-keeping incentives more closely
        with the reported safety metrics.
  \item Representative video observations: TODO
\end{itemize}

\section{Failed attempts and implementation notes}
\begin{itemize}
  \item The repository is public, so Colab cloning should not use a personal access token.
  \item Colab may include the obsolete \texttt{gym} distribution. The notebooks remove it
        before installing SafeDrive, which uses Gymnasium.
  \item Each reset seed must remain inside its configured range: validation begins at
        1000 and held-out testing begins at 3000 for corrected configurations.
  \item The initial smoke and SAC runs failed because training and callback evaluation both
        created MetaDrive's process-global engine in one Python process. Callback evaluation
        now uses a one-worker \texttt{SubprocVecEnv}.
  \item The first video attempt failed with
        \texttt{scenario\_index (seed) should be in [1000:1001)} because the vector
        environment reset without receiving seed 1000. The recorder now narrows the
        environment's valid range to the requested seed, seeds all random generators, calls
        \texttt{environment.seed(video\_seed)} immediately before reset, and loads PPO on
        its configured CPU device.
  \item Treat \texttt{/content/safedrive} as the live checkout. Mounted Drive is only for
        persistent artifacts because many small Drive-backed file operations are slow.
  \item Add dated notes and complete exception messages here; detailed stack traces remain
        in the run log files.
\end{itemize}

\section{Unresolved issues}
\begin{itemize}
  \item Repeat the already-passing 10-scenario IDM reproducibility gate once in Colab as
        a runtime smoke check before spending credits on training.
  \item Seeds 2000--2099 were consumed by the second audit. Do not report them as held out;
        corrected frozen configurations use 3000--3099 once.
  \item Run the short no-traffic diagnostic pilots and enforce their gates before another
        full run.
  \item Final numerical results and report placeholders: TODO after corrected long runs.
  \item Whether 1,000,000 timesteps materially improves either agent: optional manual test.
  \item Local macOS versus Colab runtime differences: TODO after both smoke tests.
\end{itemize}

\section{Deferred Phase-2 ideas}
Seed sweeps, richer observation studies, and explicit safety interventions are intentionally
outside Phase 1. Revisit them only after the current IDM--PPO--SAC matrix has complete,
reproducible artifacts and a finished report.

\end{document}
